\chapter{绪论}

\section{研究背景与意义}

考勤管理是学校、实验室、企业和机关单位日常管理中的基础工作。准确、及时的考勤数据能够反映人员到岗情况，为课堂管理、实验室安全、绩效统计和异常追踪提供依据。常见考勤方式包括人工签到、纸质登记、IC卡刷卡、指纹识别等。这些方式虽然实现成本较低，但在实际使用中存在明显不足：人工签到效率较低，面对集中签到时容易排队，代签行为也不易发现；纸质登记需要人工汇总和保存，数据检索、统计和追溯成本较高；IC卡依赖物理介质，容易出现忘带、遗失和代刷问题；指纹识别属于接触式采集方式，在公共卫生和设备磨损方面存在限制。

人脸识别通过摄像头采集用户面部图像，利用人脸检测、特征提取和相似度匹配完成身份确认，具有非接触、交互自然、无需携带额外介质等特点。对于考勤场景而言，人脸识别能够减少人为操作步骤，提高签到效率，同时降低代签、代刷等行为发生的可能性。随着深度学习模型轻量化和边缘计算技术的发展，人脸检测、特征提取和活体检测可以在树莓派等低成本设备上运行，使人脸识别考勤系统摆脱对高性能服务器的强依赖。

树莓派体积小、功耗低、生态成熟，具备USB、网络、存储和Linux系统环境等基础能力。将人脸识别考勤系统部署在树莓派上，可以在实验室门口、教室入口或小型办公场所构建独立终端。与集中式视频上传方案相比，端侧部署具有三点优势。首先，摄像头采集、识别推理和签到写库均在设备端完成，可以减少视频流上传带来的带宽压力。其次，人员照片和考勤数据保存在本地SQLite数据库中，适合小规模场景独立运行。最后，系统可以通过局域网或Tailscale等虚拟组网方式提供远程管理入口，在保证设备自治的同时支持管理员审核和运维。

面向真实考勤场景的人脸识别系统不能只完成一次简单的人脸比对。系统还需要处理注册资料真实性、识别实时性、重复签到、照片或屏幕翻拍攻击、远程审核和异常反馈等问题。因此，本文围绕树莓派端侧部署场景，设计并实现一套完整的人脸识别考勤系统。系统以OpenCV YuNet和SFace作为最终识别链路，结合OpenVINO被动活体检测、动作挑战、人脸质量校验、审核流和管理员API，形成从成员名单导入到考勤记录查询的端到端业务闭环。

\section{国内外研究现状}

人脸识别研究经历了从人工特征到深度学习特征的演进。早期方法通常依赖主成分分析、局部二值模式等人工特征，面对光照变化、姿态变化和遮挡场景时鲁棒性有限。随着卷积神经网络的发展，深度学习方法逐渐成为主流。Schroff等提出的FaceNet模型通过三元组损失直接学习人脸嵌入空间，使同一身份样本在特征空间中距离更近、不同身份样本距离更远\cite{schroff2015facenet}。Deng等提出的ArcFace通过加性角度间隔提升类间区分度\cite{deng2019arcface}，Wang等提出的CosFace通过大间隔余弦损失增强特征判别能力\cite{wang2018cosface}。这些研究推动了人脸特征从分类输出向嵌入向量匹配转变，也为工程系统中的特征比对提供了理论基础。

人脸检测是识别流程的前置环节。Zhang等提出的MTCNN将人脸检测和关键点定位结合起来，能够输出人脸框和关键点，广泛应用于注册质量校验和姿态估计场景\cite{zhang2016mtcnn}。近年来，端侧部署对模型体积和推理速度提出了更高要求。YuNet面向轻量化人脸检测任务设计，能够在较小模型规模下输出人脸框与关键点信息，适合CPU设备上的实时检测\cite{wu2023yunet}。SFace通过约束特征空间提升人脸特征的可分性和鲁棒性\cite{zhong2021sface}，并可与OpenCV提供的检测和识别接口结合，形成可部署的人脸识别链路\cite{opencv_dnn_face}。

活体检测是人脸识别系统安全性的重要组成部分。人脸考勤系统面对的攻击方式通常包括纸质照片、手机屏幕翻拍和预录视频等。活体检测方法可以分为主动式和被动式两类。主动式方法要求用户完成眨眼、摇头、张嘴等动作，再由系统判断动作是否真实；被动式方法不要求用户额外配合，而是利用纹理、反光、频域特征或生理信号进行判断。相关综述指出，深度学习活体检测方法需要同时关注跨设备、跨光照和跨攻击介质的泛化能力\cite{yu2023fas}。在工程实现中，主动式方法安全性较高但会增加用户操作负担，被动式方法体验更自然但对摄像头质量和现场光照更敏感。

在系统工程方面，轻量Web框架、嵌入式数据库和虚拟组网技术为端侧考勤系统提供了成熟基础。FastAPI支持类型提示和自动接口文档，适合构建设备侧HTTP服务\cite{fastapi}；SQLite以单文件形式提供关系型数据库能力，不需要额外部署数据库服务器\cite{sqlite}；Tailscale基于WireGuard构建私有网络，可让不同网络环境中的设备形成受控访问链路\cite{tailscale}。综合来看，现有人脸识别、活体检测和Web系统技术已经具备较好的可用性，但将这些技术组织成一套稳定的考勤系统，仍然需要面向注册、审核、识别、签到和管理流程进行工程化设计。

\section{本文主要工作}

本文围绕基于树莓派的人脸识别考勤系统，完成了面向普通用户、树莓派终端和管理员的主要功能设计与实现。系统主要功能用例如图\ref{fig:chapter1_work}所示。

\begin{figure}[H]
    \centering
    \resizebox{0.96\textwidth}{!}{%
    \begin{tikzpicture}[
        usecase/.style={ellipse, draw=black, align=center, minimum width=2.95cm, minimum height=0.82cm, fill=white},
        boundary/.style={rectangle, draw=black, rounded corners=4pt, inner sep=0.35cm},
        assoc/.style={-, thick},
        include/.style={->, >=stealth, dashed, thick},
        extend/.style={->, >=stealth, dashed, thick}
    ]
        \coordinate (user) at (-5.9,1.45);
        \coordinate (terminal) at (-5.9,-1.6);
        \coordinate (admin) at (5.9,0);

        \foreach \actor/\name in {user/普通用户,terminal/树莓派终端,admin/管理员} {
            \draw (\actor)+(0,0.34) circle (0.14);
            \draw (\actor)+(0,0.20) -- +(0,-0.25);
            \draw (\actor)+(-0.32,0.02) -- +(0.32,0.02);
            \draw (\actor)+(0,-0.25) -- +(-0.28,-0.60);
            \draw (\actor)+(0,-0.25) -- +(0.28,-0.60);
            \node[align=center] at ($(\actor)+(0,-0.92)$) {\name};
        }

        \node[usecase] (uc1) at (-1.8,2.15) {首次注册};
        \node[usecase] (uc2) at (-1.8,0.95) {动作挑战采集};
        \node[usecase] (uc3) at (1.65,1.55) {人脸质量校验};
        \node[usecase] (uc4) at (1.65,0.35) {人脸档案审核};
        \node[usecase] (uc5) at (-1.8,-0.95) {端侧人脸识别};
        \node[usecase] (uc6) at (1.65,-0.95) {活体检测};
        \node[usecase] (uc7) at (-1.8,-2.15) {自动签到};
        \node[usecase] (uc8) at (1.65,-2.15) {考勤记录查询};
        \node[usecase] (uc9) at (1.65,2.75) {设备状态管理};
        \node[usecase] (uc10) at (1.65,-3.25) {重新采集资料};

        \node[boundary, fit=(uc1)(uc2)(uc3)(uc4)(uc5)(uc6)(uc7)(uc8)(uc9)(uc10),
              label={[yshift=-0.2cm]above:基于树莓派的人脸识别考勤系统}] (system) {};

        \draw[assoc] (user) -- (uc1);
        \draw[assoc] (user) -- (uc2);
        \draw[assoc] (user) -- (uc10);
        \draw[assoc] (terminal) -- (uc5);
        \draw[assoc] (terminal) -- (uc7);
        \draw[assoc] (admin) -- (uc4);
        \draw[assoc] (admin) -- (uc8);
        \draw[assoc] (admin) -- (uc9);

        \draw[include] (uc2) -- node[above]{\scriptsize $\langle\langle$include$\rangle\rangle$} (uc3);
        \draw[include] (uc5) -- node[above]{\scriptsize $\langle\langle$include$\rangle\rangle$} (uc6);
        \draw[include] (uc7) -- node[left]{\scriptsize $\langle\langle$include$\rangle\rangle$} (uc5);
        \draw[extend] (uc10) -- node[right]{\scriptsize $\langle\langle$extend$\rangle\rangle$} (uc2);
    \end{tikzpicture}}
    \caption{系统主要功能用例图}
    \label{fig:chapter1_work}
\end{figure}

图\ref{fig:chapter1_work}从UML用例视角概括了本文系统的主要功能边界。普通用户主要参与首次注册、动作挑战、人脸采集和驳回后的重新采集；树莓派终端主要承担端侧人脸识别、活体检测和自动签到；管理员主要负责人脸档案审核、考勤记录查询和设备状态管理。图中实线表示参与者与用例之间的关联关系，虚线箭头表示用例之间的包含或扩展关系。与普通流程图相比，用例图更能体现本文工作的系统边界、参与对象和功能完整性。

第一，设计并实现完整的考勤业务链路。系统从本地CSV成员名单开始，支持成员预置、用户首次注册、密码修改、人脸采集、管理员审核、终端实时识别、自动签到和考勤记录查询，覆盖小规模考勤系统的主要业务流程。

第二，设计并实现注册端动作挑战和质量校验机制。系统采用正脸、侧转、再正脸的三步动作挑战，后端保存挑战会话并冻结最终可信注册照，避免前端通过挑战后更换图片。系统实现单人检测、人脸尺寸、模糊度、亮度、关键点完整性、姿态和翻拍风险等质量校验，提高入库资料质量。

第三，设计并实现基于YuNet和SFace的终端识别链路。系统使用YuNet进行人脸检测，使用SFace提取人脸特征向量，通过余弦相似度与已审核人脸库进行身份匹配，并结合单人入镜、人脸面积、活体检测、考勤时段和重复签到冷却等门禁保证签到可靠性。

第四，设计并实现被动活体检测和屏幕重放风险分析。系统通过OpenVINO加载anti-spoof-mn3模型，对裁剪后的人脸图像输出真人概率，并使用2秒滑动窗口融合多帧结果。同时，系统结合反光、边缘密度、频域高频能量和矩形边框等轻量信号判断屏幕重放风险。

第五，设计并实现远程管理接口。系统在树莓派端提供\texttt{/api/admin/*}管理员接口，支持设备状态、成员名单、用户、人脸档案审核、考勤记录和快照访问等功能。管理员端可以通过Tailscale网络访问树莓派接口，实现远程审核和运维。

\section{论文组织结构}

本文正文围绕基于树莓派的人脸识别考勤系统展开，共分为五章。

第一章为绪论，介绍课题背景与意义、国内外研究现状、本文主要工作和论文组织结构。

第二章为毕业设计相关关键技术，介绍系统所依赖的OpenCV YuNet、SFace、OpenVINO活体检测、FastAPI、Jinja2、Bootstrap、SQLite和Tailscale等技术。

第三章为系统需求分析，分析系统参与者、功能需求、非功能需求和核心业务流程，为后续系统设计提供依据。

第四章为系统设计，围绕系统类毕业设计展开，介绍功能结构、总体架构、部署方式、注册审核设计、端侧识别设计、数据模型和接口组织方式。

第五章为系统实现与测试分析，说明用户门户、注册采集、终端识别、管理员接口和数据持久化的实现过程，并给出功能测试结果和性能分析。

最后为结论，总结全文工作，分析系统不足，并提出后续改进方向。
